\chapter{مواد لا تحدد نوع العصبون}
\label{sec:othermediators}

في الفصل~\ref{sec:neurontypes} فرزنا العصبونات بحسب ناقلها العصبي الرئيسي وحصلنا على عشرة أنواع (الجدول~\ref{tab:types}). وقد حان الوقت لإضافة التعقيدات. فالعصبون لا يقذف الناقل الرئيسي وحده، وفي الدماغ تعمل إلى جانب النواقل الرئيسية مواد أخرى. وهي تغيّر طريقة الشبكة في نقل الإشارات وتخزينها.

وإلى جانب ناقل العناوين وناقل البثّ (القسم~\ref{sec:buses}) هناك ناقل ثالث: \emph{القناة العكسية}. فبعض المواد لا يقذفها المرسل، بل \emph{المستقبِل}، فتعبر الشق راجعة إلى مخرج المرسل، وتغيّر كمية الناقل التي سيقذفها في المرة التالية. وهذا يشبه في شبكات الاتصال إشعار الاستلام الذي يستعمله المرسل لضبط إرساله. وبهذه القناة بالذات تستعين العصبونات حين تغيّر أوزان روابطها: فعبرها يعرف مخرج المرسل فعالية المستقبِل، أي العامل الثاني في قواعد التعلم في الفصل~\ref{sec:dofa}.

القائمة الكاملة للنواقل المعروفة طويلة: أكثر من مئة مادة، أغلبها سلاسل بروتينية قصيرة هي الببتيدات العصبية. لكنها كلها تندرج في بضعة أصناف، والمواد داخل كل صنف تتصرف تصرفًا متشابهًا. وقد جُمعت المواد التي لا تكون ناقلًا رئيسيًا في الجدول~\ref{tab:othermediators}، مع الأسئلة الثلاثة نفسها التي يطرحها المهندس: على أي ناقل تسير الإشارة، وهل تعمل بسرعة، وما علامتها المعتادة. وهي لا تحدد نوع العصبون: فإما أن تُقذف مع الناقل الرئيسي، وإما أن تقذفها خلايا غير عصبونية، وإما أن تسير في الاتجاه المعاكس.

\begin{table}[t]
\centering
\footnotesize
\setlength{\tabcolsep}{4pt}
\renewcommand{\arraystretch}{1.12}
\begin{tabular}{>{\raggedright}p{2.25cm}>{\raggedright}p{2.8cm}>{\raggedright}p{1.8cm}>{\raggedright}p{1.5cm}>{\centering}p{0.8cm}>{\raggedright\arraybackslash}p{4.3cm}}
\toprule
الصنف & المادة & الناقل & السرعة & العلامة & الدور في الحساب \\
\midrule
الأحماض الأمينية & دي-سيرين & موضعي & بطيء & و & مفتاح ثانٍ لمستقبل الغلوتامات: شرط ”و“ \\
\midrule
الأمينات الأحادية & الأمينات الأثرية & بثّ & بطيء & $\pm$ & ضبط دقيق لقنوات الأمينات الأحادية \\
\midrule
\multirow{2}{2.25cm}{البيورينات}
 & \foreignlanguage{english}{ATP} & عنواني & سريع وبطيء & $+$ & ناقل مرافق؛ إشارة بين العصبونات والدِّبق \\
 & الأدينوسين & حجمي & بطيء & $-$ & يتراكم مع الفعالية: عدّاد للحِمل~\cite{Dunwiddie2001} \\
\midrule
\multirow{8}{2.25cm}{الببتيدات العصبية (أكثر من مئة)}
 & الإنكيفالينات، الإندورفينات، الداينورفين & حجمي & بطيء & $-$ & كبت النقل \\
 & المادة \foreignlanguage{english}{P} & حجمي & بطيء & $+$ & تعزيز بطيء \\
 & الببتيد العصبي \foreignlanguage{english}{Y} & حجمي & بطيء & $-$ & تثبيط بطيء \\
 & السوماتوستاتين & حجمي & بطيء & $-$ & تثبيط بطيء، يرافق الغابا \\
 & الببتيد المعوي الفعال في الأوعية (\foreignlanguage{english}{VIP}) & حجمي & بطيء & $+$ & رفع التثبيط: تثبيط العصبونات المثبِّطة \\
 & الكوليسيستوكينين & حجمي & بطيء & $\pm$ & يرافق الغابا في بعض العصبونات المثبِّطة \\
 & الأوركسين & بثّ & بطيء & $+$ & مثبِّت لنظام اليقظة \\
 & الأوكسيتوسين، الفازوبريسين & بثّ & بطيء & $\pm$ & ضبط شامل للسلوك \\
\midrule
\multirow{2}{2.25cm}{الغازات}
 & أكسيد النيتريك \foreignlanguage{english}{NO} & عكسي، حجمي & بطيء & $\pm$ & يعبر الأغشية؛ إشارة عكسية عند تغيّر الأوزان~\cite{Garthwaite2008} \\
 & \foreignlanguage{english}{CO}، \foreignlanguage{english}{H$_2$S} & حجمي & بطيء & $\pm$ & دوره قليل الدراسة \\
\midrule
الدهون & الكانابينويدات الداخلية (الأنانداميد، \foreignlanguage{english}{2-AG}) & عكسي & بطيء & $-$ & المستقبِل يقلّل القذف عند المرسل: تغذية راجعة على الوزن~\cite{WilsonNicoll2001} \\
\bottomrule
\end{tabular}
\caption{مواد لا تحدد نوع العصبون. \emph{الناقل} و\emph{السرعة} و\emph{العلامة} كما في الجدول~\ref{tab:types}؛ ويضاف إليها: الناقل الحجمي، حيث تنتشر المادة في الحيز بين الخلايا حول مكان القذف؛ والعكسي، حيث تسير من المستقبِل راجعة إلى المرسل؛ والموضعي، حيث تعمل قرب مكان القذف. ”و“ إشارة إذن: لا تولّد تيارًا بذاتها، لكن من دونها لا يعمل مدخل آخر، كالمدخل الثاني لبوابة ”و“ (\foreignlanguage{english}{AND}) المنطقية. والببتيدات العصبية تُقذف دائمًا تقريبًا مع الناقل الرئيسي، وفي الغالب عند تردد نبضات مرتفع~\cite{vandenPol2012}.}
\label{tab:othermediators}
\end{table}
